\lecture{2}{Tue 30 Aug 2022 19:05}{Groups}

\begin{corollary}[Euclid]
	There are infinitely many primes.
\end{corollary}

\begin{proposition}
	Let $p_1 = 2$, $p_2 = 3, \ldots$ be a enumeration of primes. Then the smallest divisor of \[
		\left( p_1\ldots p_n \right) ^{\left( p_1 \ldots p_n \right) ^{\left( p_1 \ldots p_n \right) }}
	\] is $p_{n+1}$.
\end{proposition}

\section{Groups}
\subsection{Definition and Examples}
\begin{definition}
	A \textit{binary operation} $\psi$ over a set $G$ is a function \[
	\psi : G \times  G \to  G
	.\] 
\end{definition}

\begin{definition}[Group]
	A \textit{group} is a set $G$ equipped with a binary operation $\psi$ such that 
	\begin{enumerate}
		\item (G0)(closure) $\psi$ is closed in $G$,
		\item (G1)(associativity) $\psi$ is associative,
		\item (G2)(identity) Exists an identity $e \in G$ such that $ae = ea = a$ for $\forall a \in G$, and
		\item (G3)(inverse) For $\forall a \in G$, there exists $a^{-1} \in G$ such that $a a^{-1} = a^{-1} a = e$.
	\end{enumerate}
\end{definition}

\begin{example}
	The \textit{trivial group} $\{e\} $.
\end{example}

\begin{example}
	The \textit{symmetry group} $A(S)$ with \textit{composition} $\circ$ as the binary operation.
\end{example}

\begin{example}
	$(\mathbb{Z}, +)$, $\left( \mathbb{R}^\times , \times  \right) $ are addition group over integers and multiplication group over nonzero reals.
\end{example}

\begin{example}
	$\left( \mathbb{Z}_n, +_n \right)$ where $\mathbb{Z}_n = \{0,1,\ldots,n -1 \} $ is the group of addition mod $n$.
\end{example}

\begin{example}
	$(\text{SL}_2(\mathbb{R}), *)$ is the \textit{special linear group} of $2\times 2$ matrices with determinant 1. This is precisely the positive isometries.
\end{example}

\begin{example}
	$(D_{2n}, *)$ is the \textit{dihedral group} of symmetry of regular $n$-sided polygons.
	
	This group can be constructed explicitly by the $n$ th roots of unity. The group elements are then 
	\[
	\sigma : e^{2 \pi i \frac{a}{n}} \mapsto e^{2 \pi i \frac{a+1}{n}}
	\] the counterclockwise rotation by $e^{\frac{2}{n} i \pi}$, and 
	\[
	\tau: e^{2 \pi i \frac{a}{n}} \mapsto e^{-2 \pi i \frac{a}{n}}
	\] the reflection over real axis.

	The mappings $1, \sigma, \sigma^2, \ldots, \sigma^{n-1}$ are distinct. The mappings $\tau, \tau\sigma, \ldots, \tau\sigma^{-1}$ are all distinct and distinct from $1, \sigma, \ldots, \sigma^{-1}$.

	Check: $\sigma \tau = \tau \sigma^{-1} $. \[
	\sigma\tau\left( e^{2 \pi \frac{a}{n}} \right)  = \sigma\left( e^{- 2 \pi i \frac{a}{n}} \right) = e^{-2 \pi i \frac{1-a}{n}}
	.\] 
	\[
	\tau\sigma^{-1}(e^{2 \pi i \frac{a}{n}} = \tau\left( e^{2 \pi i \frac{a^{-1} }{n}} \right) =e^{2 \pi i \frac{a -1}{n}}
	.\] 
\end{example}

\begin{remark}
	Geometers often refer to the group $D_{2n}$ as $D_n$.
\end{remark}

\begin{definition}
	A group $(G, *)$ is said to be a finite group if $\text{card }G < \infty $, in this case we say the \textit{order} of group $G$ is $|G| = \text{card } G$.
\end{definition}

\begin{example}
	\begin{enumerate}
\item $|A(S)| = n!$ where $|S| = n$
\item $|\mathbb{Z}| = +\infty $
\item $|\mathbb{Z}_n| = n$
\item $|\mathbb{R}^\times | = +\infty $
\item $\left| \text{SL}_2(\mathbb{R}) \right| = +\infty $ 
\item (harder) one can consider $\text{SL}_2 \left( \mathbb{Z}_p \right) $ has cardinality $(p -1) p (p +1)$
\item $|D_{2n}| = 2n$
	\end{enumerate}
\end{example}

\begin{definition}
	A group $(G, *)$ is called \textit{abelian} if $*$ is commutative, i.e.,  for all $a,b \in G$, we have $a*b = b*a$.
\end{definition}
\begin{example}
	\begin{enumerate}
		\item In general $A(S)$ is non-abelian
		\item $(\mathbb{Z},+)$ is abelian
		\item $(\mathbb{Z}_n,+)$ is abelian
		\item $\left( \mathbb{R}^\times  ,\times  \right) $ is abelian
		\item $\text{SL}_2 (\mathbb{R}) $ is non-abelian (matrices do not commute in general)
		\item $D_{2n}$ is non-abelian when $n>1$
	\end{enumerate}
\end{example}

\subsection{The most basic properties}

\begin{lemma}[2.2.1]
	Suppose that $G$ is a group, then
	\begin{enumerate}
		\item the identity element $e \in G$ is unique;
		\item for every $a \in G$, the inverse $a^{-1} \in G$ is unique;
		\item for all $\left( a^{-1}  \right)^{-1}=a $;
		\item for all $a,b \in G$, $\left( ab  \right) ^{-1} = b^{-1} a^{-1} $.
	\end{enumerate}
\end{lemma}
\begin{proof}
\begin{enumerate}
	\item Suppose $e_1, e_2 \in G$ satisfy identity axiom. Then \[
	e_1e_2 = e_1 = e_2
	\] where we applied the identity property twice. 
\item Suppose $b_1, b_2 \in G$ are both inverses for $a \in G$. Then if $e$ is the identity of $G$,
	\begin{align*}
		b_1ab_2 &= b_1(ab_2) = b_1e=b_1 \\
		&= (b_1a)b_2=eb_2=b_2
		.
	\end{align*}
\item \[
a^{-1} a = e = aa^{-1} 
.\] So $a$ is the inverse of $a^{-1}$.
\item Let $c = \left( ab  \right) ^{-1}$. Then $(ab)c = e$. Then $b^{-1} a^{-1} (ab)c = b^{-1} a^{-1} $. Use associativity, inverse and identity twice, we simplify the LHS to $c$.
\end{enumerate}
\end{proof}

\begin{lemma}[cancellation rule]
	Suppose  $G$ is a group, and $a,b,c \in G$. Then
	\begin{enumerate}
		\item whenever $ab=ac$, one has $b=c$;
		\item whenever $ba=ca$, one has $b=c$.
	\end{enumerate}
\end{lemma}
\begin{proof}
\begin{enumerate}
	\item Since $ab=ac$, $a^{-1} ab = a^{-1} ac$. Use associativity, inverse and identity axiom, and we simplify to $b=c$.
	\item Similar.
\end{enumerate}
\end{proof}

\subsection{Subgroups}

\begin{definition}[subgroups]
	Suppose $H$ is a non-empty subset of a group $G$ equipped with a binary operation. If $H$ forms a group with respect to the same binary operation, then we call $H$ a \textit{subgroup} of $G$ and write $H\le G$.
\end{definition}

\begin{observation}
	Suppose $G$ is a group with identity element $e$, then there are two \textit{trivial subgroups}
	\begin{itemize}
		\item $G\le G$ 
		\item $\{\} \le G$
	\end{itemize}
	Any subgoup $H$ of $G$ that is nontrivial is called a \textit{proper subgroup}.
\end{observation}

\begin{lemma}[subgroup criterion]
	Suppose that $G$ is a group, and $\emptyset \neq H \subset G$. Then $H\le G$ iff for all $a,b \in H$ one has $ab^{-1} \in H$.
\end{lemma}
\begin{proof}
	If  $H\le G$ then for all $b \in H$, one has $b^{-1}  \in H$, so by closure we have for all $a,b \in H$, $ab^{-1}  \in H$.

	Let $a \in H$ with $\emptyset \neq H \subset G$, we may assume that for $b \in H$ we have $ab^{-1} \in H$. So in particular $aa^{-1} = e \in H$. Also, since $e \in H$, we have $ea^{-1} \in H$, so $a^{-1}  \in H$. For associativity, this is inherited from $G$. Finally, to check closure, assume that whenever $a,b \in H$, $b^{-1}  \in H$, so $ab = a(b^{-1} )^{-1} \in H$.
\end{proof}

\begin{example}
	\begin{enumerate}
	\item Let $n \in \mathbb{N}$ (think $n\ge 2$ ). Consider $n\mathbb{Z}$ as a subset of $\mathbb{Z}$ with $(\mathbb{Z},+)$ a group. Then $n\mathbb{Z}\le \mathbb{Z}$. Proof: If $a, b \in n\mathbb{Z}$, then $n|a$ and $n | b$, so $n|(a-b)$, where $a-b \in n\mathbb{Z}$. So subgroup criterion applies.
	\item Let $\left( \mathbb{C}^\times ,\times  \right) $ be the group of non-zero complex numbers with ordinary multiplication. Consider  \[
	\Pi = \{z \in \mathbb{C}: |z|=1\}
	.\] 
	Check that $\Pi \le  \mathbb{C}^\times $.
\item Consider $\mu_n = \{e^{2 \pi i \frac{a}{n}}: 0\le a<n, a\in \mathbb{Z}\} 	$. Now \[
\mu_n \le \Pi \le  \mathbb{C}^{\times }
.\] 
\item Let $(G, \cdot)$ be a group and $a \in G$. Then $\left\langle x \right\rangle = \{a^n: n \in \mathbb{Z}\} $ forms a subgroup of $G$.
	\begin{definition}[cyclic group]
		A group G generated by a single element $a$ so that $G = \left\langle a \right\rangle $ is called \textit{cyclic}.
	\end{definition}
\item We have $\left\langle 2 \right\rangle \le \mathbb{Z}_6 = \left\langle 1 \right\rangle $. $\left\langle 3 \right\rangle \le \mathbb{Z}_6$. $\mathbb{Z}_7$ does not have any proper subgroups.
\item (Centralizer) When $a \in G$, define \[
C(a) := \{g \in G: ga = ag\} 
.\] 
\begin{theorem}
	For each $a \in G$, one has $C(a) \le G$.
\end{theorem}
\begin{proof}
	Observe that whenever $h \in C(a)$, then
	\begin{align*}
		ha&=ah\\
		a = h^{-1}  h a &= h^{-1}  a h \\
		ah^{-1} &= h^{-1} ahh^{-1} =h^{-1} a \\
	\end{align*}
	. Hence $h^{-1} ha = h_{-1} ah$, where $h^{-1}  \in C(a)$. Thus, whenever $g,h \in C(a)$, one has \[
	a(gh^{-1} )=(ag)h^{-1} =(ga)h^{-1} =g(ah^{-1} )=g(h^{-1} a)=(gh^{-1} )a
	.\] 
	So $gh^{-1} \in C(a)$. Subgoup criterion applies.
\end{proof}
\item (Center) \[
Z(G)= \{z\in G: zx=xz \text{ for all }x \in G\}
.\] 
$Z(G)$ is a subgroup of $G$.
\begin{proof}
Observe that $Z(G) = \bigcap_{a \in G} C(a)$. Note: if $K$ and $H$ are both subgroups of $G$, then $K \cap H\le G$. (In case of uncountable groups we need axiom of choice to settle this down.)
\end{proof}


	\end{enumerate}
\end{example}
